182
«من بزرگترین هنرمند رابطه‌ی جهان هستم»، او به سمت من غرغرکنان گفت.
«چرا دوست‌دختر ندارم؟»
«خب، شاید چون بزرگترین هنرمند رابطه‌ی جهان هستی.»

پس از یک مدت طولانی سکوت، میستری از کارولین خواست او را به باشگاه رقاصی ببرد که دوست‌دختر سابقش پاتریشیا در آن کار می‌کرد. او را در پارکینگ پیاده کرد و سپس مرا برای گذراندن شب به خانه‌ای در حومه شهر برد که با مادر، خواهر و برادرش در آن زندگی می‌کرد. در حالی که داشتم با کارولین نقش پدر را بازی می‌کردم، میستری در حالت سقوط بود. این اولین بار بود که خانواده‌اش را می‌دیدم. پیاده کردن او در باشگاه کار اشتباهی بود. دیدن پاتریشیا او را به هم ریخته بود. نه تنها او حاضر نبود به او برگردد، بلکه به او گفته بود که با افراد دیگری ملاقات کرده است.

مادرش در آستانه در با گریه نوزادی در آغوش به ما خوش‌آمد گفت — نوزاد دوست‌دختر نوجوانم.

«می‌خوای او را بغل کنی؟» کارولین پرسید. فکر می‌کنم واکنش کلیشه‌ای این است که بگوییم ترسیده بودم و واقعیت بر من غلبه کرد و می‌خواستم از آنجا بیرون بروم.

اما اینطور نبود. می‌خواستم او را بغل کنم. یک‌جور جالب بود. من برای داشتن چنین ماجراجویی‌هایی وارد این کار شدم، برای اینکه اولین بار نوزادی را در آغوش بگیرم و بپرسم: «مادرش چه انتظاری از من دارد؟»

«سه ساعت در روز تمرین کرده.» کارولین از طریق تلفن گفت. «پونزده پوند وزن کم کرده و بدنش از صد بهتر به نظر می‌رسه. کارهایی که یک دختر وقتی عصبی می‌شود انجام می‌دهد. لعنتی.»

«فکر نکن چقدر خوب به نظر می‌آید.» به او توصیه کردم. «به عیب‌ها فکر کن و در ذهنت بزرگشان کن. این کار را آسان‌تر می‌کند.»

«می‌دانم که از نظر عقلی، اما احساسی گند زده‌ام. احساس می‌کنم دارم روی خاکستر داغ کشیده می‌شوم. همه چیز وقتی دوباره او را دیدم بر سرم خراب شد. آن بدن جذاب، خط‌های برنزه. او جذابترین رقاص کلوب بود. و من نمی‌توانم او را داشته باشم. کری با دوست‌پسرش برگشته. و من از تلاش برای زیبا کردن محل جدیدم خسته‌ام. برای چی؟»

«رفیق، تو هنرمند رابطه‌ای. صدها نفر مثل پاتریشیا در آنجا وجود دارند. و تو می‌توانی در یک شب آنها را پیدا کنی.»

«من هنرمند رابطه نیستم. من عاشقم. من زنان را دوست دارم. قسم می‌خورم، حتی به سه‌نفری هم فکر نمی‌کنم. حالا با رفتن با پاتریشیا خواهم خوشحال شدم. دائماً به او فکر می‌کنم. هر دقیقه روز دلم براش تنگ می‌شه.»

میستری تا زمانی که پاتریشیا او را رد نکرده بود به او فکر نکرده و با او صحبت هم نکرده بود. حالا او تبدیل به وسواسی شده بود. نظریه‌های خودش درباره جذب دوباره به او ضربه زده بود. پاتریشیا او را کنار گذاشته بود. اما برای او این یک تکنیک نبود — واقعیت داشت.

به عنوان یک شعبده‌باز که به بهره‌گیری از ساده‌لوحی دیگران عادت داشت، میستری هیچ صبری برای چیزهایی روحانی یا ماوراءالطبیعه نداشت. دین او داروین بود. عشق، برای او، صرفاً یک تکانه تکاملی بود که به انسان‌ها اجازه می‌داد دو هدف اصلی را برآورده کنند: بقا و تولیدمثل. او آن تکانه را «پایر باندینگ» می‌نامید.

«عجیبه که چقدر پایر باندینگ قویه»، او گفت. «حالا احساس تنهایی می‌کنم.»

«بهت می‌گم چی. فردا می‌آییم دنبالت و در حومه شهر با ما بازی می‌کنی. خوشحالت می‌کنه.»